% Options for packages loaded elsewhere
\PassOptionsToPackage{unicode}{hyperref}
\PassOptionsToPackage{hyphens}{url}
\PassOptionsToPackage{dvipsnames,svgnames,x11names}{xcolor}
\documentclass[
  american,
  11pt,
]{article}
\usepackage{xcolor}
\usepackage[margin=1in]{geometry}
\usepackage{amsmath,amssymb}
\setcounter{secnumdepth}{5}
\usepackage{iftex}
\ifPDFTeX
  \usepackage[T1]{fontenc}
  \usepackage[utf8]{inputenc}
  \usepackage{textcomp} % provide euro and other symbols
\else % if luatex or xetex
  \usepackage{unicode-math} % this also loads fontspec
  \defaultfontfeatures{Scale=MatchLowercase}
  \defaultfontfeatures[\rmfamily]{Ligatures=TeX,Scale=1}
\fi
\usepackage{lmodern}
\ifPDFTeX\else
  % xetex/luatex font selection
\fi
% Use upquote if available, for straight quotes in verbatim environments
\IfFileExists{upquote.sty}{\usepackage{upquote}}{}
\IfFileExists{microtype.sty}{% use microtype if available
  \usepackage[]{microtype}
  \UseMicrotypeSet[protrusion]{basicmath} % disable protrusion for tt fonts
}{}
\makeatletter
\@ifundefined{KOMAClassName}{% if non-KOMA class
  \IfFileExists{parskip.sty}{%
    \usepackage{parskip}
  }{% else
    \setlength{\parindent}{0pt}
    \setlength{\parskip}{6pt plus 2pt minus 1pt}}
}{% if KOMA class
  \KOMAoptions{parskip=half}}
\makeatother
\usepackage{color}
\usepackage{fancyvrb}
\newcommand{\VerbBar}{|}
\newcommand{\VERB}{\Verb[commandchars=\\\{\}]}
\DefineVerbatimEnvironment{Highlighting}{Verbatim}{commandchars=\\\{\}}
% Add ',fontsize=\small' for more characters per line
\newenvironment{Shaded}{}{}
\newcommand{\AlertTok}[1]{\textcolor[rgb]{1.00,0.00,0.00}{\textbf{#1}}}
\newcommand{\AnnotationTok}[1]{\textcolor[rgb]{0.38,0.63,0.69}{\textbf{\textit{#1}}}}
\newcommand{\AttributeTok}[1]{\textcolor[rgb]{0.49,0.56,0.16}{#1}}
\newcommand{\BaseNTok}[1]{\textcolor[rgb]{0.25,0.63,0.44}{#1}}
\newcommand{\BuiltInTok}[1]{\textcolor[rgb]{0.00,0.50,0.00}{#1}}
\newcommand{\CharTok}[1]{\textcolor[rgb]{0.25,0.44,0.63}{#1}}
\newcommand{\CommentTok}[1]{\textcolor[rgb]{0.38,0.63,0.69}{\textit{#1}}}
\newcommand{\CommentVarTok}[1]{\textcolor[rgb]{0.38,0.63,0.69}{\textbf{\textit{#1}}}}
\newcommand{\ConstantTok}[1]{\textcolor[rgb]{0.53,0.00,0.00}{#1}}
\newcommand{\ControlFlowTok}[1]{\textcolor[rgb]{0.00,0.44,0.13}{\textbf{#1}}}
\newcommand{\DataTypeTok}[1]{\textcolor[rgb]{0.56,0.13,0.00}{#1}}
\newcommand{\DecValTok}[1]{\textcolor[rgb]{0.25,0.63,0.44}{#1}}
\newcommand{\DocumentationTok}[1]{\textcolor[rgb]{0.73,0.13,0.13}{\textit{#1}}}
\newcommand{\ErrorTok}[1]{\textcolor[rgb]{1.00,0.00,0.00}{\textbf{#1}}}
\newcommand{\ExtensionTok}[1]{#1}
\newcommand{\FloatTok}[1]{\textcolor[rgb]{0.25,0.63,0.44}{#1}}
\newcommand{\FunctionTok}[1]{\textcolor[rgb]{0.02,0.16,0.49}{#1}}
\newcommand{\ImportTok}[1]{\textcolor[rgb]{0.00,0.50,0.00}{\textbf{#1}}}
\newcommand{\InformationTok}[1]{\textcolor[rgb]{0.38,0.63,0.69}{\textbf{\textit{#1}}}}
\newcommand{\KeywordTok}[1]{\textcolor[rgb]{0.00,0.44,0.13}{\textbf{#1}}}
\newcommand{\NormalTok}[1]{#1}
\newcommand{\OperatorTok}[1]{\textcolor[rgb]{0.40,0.40,0.40}{#1}}
\newcommand{\OtherTok}[1]{\textcolor[rgb]{0.00,0.44,0.13}{#1}}
\newcommand{\PreprocessorTok}[1]{\textcolor[rgb]{0.74,0.48,0.00}{#1}}
\newcommand{\RegionMarkerTok}[1]{#1}
\newcommand{\SpecialCharTok}[1]{\textcolor[rgb]{0.25,0.44,0.63}{#1}}
\newcommand{\SpecialStringTok}[1]{\textcolor[rgb]{0.73,0.40,0.53}{#1}}
\newcommand{\StringTok}[1]{\textcolor[rgb]{0.25,0.44,0.63}{#1}}
\newcommand{\VariableTok}[1]{\textcolor[rgb]{0.10,0.09,0.49}{#1}}
\newcommand{\VerbatimStringTok}[1]{\textcolor[rgb]{0.25,0.44,0.63}{#1}}
\newcommand{\WarningTok}[1]{\textcolor[rgb]{0.38,0.63,0.69}{\textbf{\textit{#1}}}}
\usepackage{longtable,booktabs,array}
\usepackage{caption}
\captionsetup[table]{skip=6pt}
\newcounter{none} % for unnumbered tables
\usepackage{calc} % for calculating minipage widths
% Correct order of tables after \paragraph or \subparagraph
\usepackage{etoolbox}
\makeatletter
\patchcmd\longtable{\par}{\if@noskipsec\mbox{}\fi\par}{}{}
\makeatother
% Allow footnotes in longtable head/foot
\IfFileExists{footnotehyper.sty}{\usepackage{footnotehyper}}{\usepackage{footnote}}
\makesavenoteenv{longtable}
\ifLuaTeX
\usepackage[bidi=basic,shorthands=off]{babel}
\else
\usepackage[bidi=default,shorthands=off]{babel}
\fi
\ifLuaTeX
  \usepackage{selnolig} % disable illegal ligatures
\fi
\setlength{\emergencystretch}{3em} % prevent overfull lines
\providecommand{\tightlist}{%
  \setlength{\itemsep}{0pt}\setlength{\parskip}{0pt}}
\usepackage{seqsplit}
\usepackage{bookmark}
\IfFileExists{xurl.sty}{\usepackage{xurl}}{} % add URL line breaks if available
\urlstyle{same}
% fallback for those not using the hyperref driver hyperxmp:
\makeatletter
\@ifundefined{xmpquote}{\newcommand{\xmpquote}[1]{#1}}{}
\makeatother
\hypersetup{
  pdftitle={Logarithmic influence bounds and variance counterexamples for Boolean decision-tree covariance},
  pdfauthor={Dev D. Goyal},
  pdflang={en-US},
  colorlinks=true,
  linkcolor={blue},
  filecolor={Maroon},
  citecolor={Blue},
  urlcolor={blue},
  pdfcreator={LaTeX via pandoc}}

\title{Logarithmic influence bounds and variance counterexamples for
Boolean decision-tree covariance}
\author{Dev D. Goyal}
\date{27 October 2026}

\begin{document}
\maketitle

\begin{center}
\small ORCID: \href{https://orcid.org/0009-0009-9132-744X}{0009-0009-9132-744X}
\end{center}

\section*{Abstract}\label{abstract}
\addcontentsline{toc}{section}{Abstract}

For a Boolean decision tree \(T\), let \(k_i(T)\) count the internal
nodes querying coordinate \(i\), and let \(I_i(f)\) be its coordinate
influence under uniform independent bits. We prove the weighted bound \[
\sum_{v\in\operatorname{Int}(T)}2^{-d(v)}
 \left|\operatorname{Cov}(f_{v,+},f_{v,-})\right|
\leq \sum_i I_i(f)\log_2\max\{1,k_i(T)\}.
\tag{A}
\]

The standard assumption is that no coordinate is queried twice on a
single root-to-leaf path. In particular, the signed tree covariance of
Wan, Wright and Wu is at most \(I(f)\log_2 k\) for a read-\(k\) tree.
The logarithmic factor has the exact coefficient one along the family
\(k=2^d\): the dummy-prefix construction of Wan, Wright and Wu, with a
dictator at its base, attains equality.

Shalev proved the earlier covariance bound \(2\sqrt{k}I(f)\). The
inequality above replaces the square-root dependence on the read number
by a logarithm, with an optimal coefficient for this covariance
quantity. The logarithmic potential itself was previously suggested by
Shalev; the contribution is a proof of the proposed inequality through a
minimum-of-influences estimate.

We give a four-variable read-three tree computing the two-variable AND
function with tree covariance \(5/4\) and variance \(3/4\). A second,
nine-variable witness has every query relevant and also violates the
proposed bound. Since \(\log_2 3<5/3\), this refutes the proposed
universal bound \(\mathcal C(T)\leq\operatorname{Var}(f)\log_2 k\). The
four-variable tree uses the original convention, which permits dummy
queries. The unrestricted Fourier entropy-influence conjecture remains
unresolved by this work.

\section{Definitions and scope}\label{definitions-and-scope}

A coordinate influence is the probability that flipping that input bit
changes the output. All expectations below are finite uniform averages.
For \(f:\{-1,1\}^n\to\{-1,1\}\), write \(x^{\oplus i}\) for \(x\) with
coordinate \(i\) flipped and define \[
I_i(f)=2^{-n}\#\{x:f(x)\ne f(x^{\oplus i})\},
\qquad I(f)=\sum_{i=1}^n I_i(f).
\]

For functions \(g,h\) on the same cube, \[
\operatorname{Cov}(g,h)=\mathbb E[gh]-(\mathbb E g)(\mathbb E h).
\]

A decision tree is a finite binary tree with signed Boolean leaves and
one coordinate label at each internal node. On an input \(x\), the
branch at a node labelled \(i\) is chosen by \(x_i\). We assume no
variable is queried twice along a root-to-leaf path. This is the same
convention used in the decision-tree covariance literature.

For an internal node \(v\), \(f_{v,+}\) and \(f_{v,-}\) denote the two
child functions, viewed as functions on the same full input cube by
ignoring coordinates the respective subtree does not query. Let \(d(v)\)
be its depth, with the root at depth zero. Define signed and absolute
tree covariances by \[
\mathcal C(T)=\sum_v2^{-d(v)}\operatorname{Cov}(f_{v,+},f_{v,-}),
\qquad
\mathcal A(T)=\sum_v2^{-d(v)}|\operatorname{Cov}(f_{v,+},f_{v,-})|.
\]

Both are zero for a leaf. If the two root subtrees \(T_+,T_-\) compute
\(g,h\), respectively, then \[
\mathcal C(T)=\operatorname{Cov}(g,h)+\frac{\mathcal C(T_+)+\mathcal C(T_-)}2,
\]

and the same recurrence with the absolute value of the root covariance
holds for \(\mathcal A\). In particular,
\(\mathcal C(T)\leq\mathcal A(T)\).

Let \(k_i(T)\) be the total number of internal nodes labelled \(i\) in
the entire tree. This is a global node count, not a count along one
path. Put \[
w(r)=\log_2\max\{1,r\}\quad(r\in\mathbb N),
\qquad
P(T)=\sum_i I_i(f)w(k_i(T)).
\]

The convention \(w(0)=w(1)=0\) avoids an undefined logarithm;
\(k_i(T)=0\) also implies \(I_i(f)=0\). A tree is read-\(k\) when
\(k_i(T)\leq k\) for every \(i\), where \(k\geq1\).

\section{A Boolean covariance
inequality}\label{a-boolean-covariance-inequality}

\subsection*{Lemma 1}\label{lemma-1}
\addcontentsline{toc}{subsection}{Lemma 1}

For every two Boolean functions \(g,h:\{-1,1\}^n\to\{-1,1\}\), \[
|\operatorname{Cov}(g,h)|\leq
\sum_{i=1}^n\min\{I_i(g),I_i(h)\}.
\tag{1}
\]

We first prove the covariance estimate used in the tree induction.

\textbf{Proof.} Induct on \(n\). When \(n=0\), the functions are
constant, so both sides are zero.

For \(n\geq1\), split the last coordinate and write \(g_+,g_-,h_+,h_-\)
for the restrictions to the \((n-1)\)-cube. Put \[
d_g=\frac{\mathbb E g_+-\mathbb E g_-}{2},
\qquad
d_h=\frac{\mathbb E h_+-\mathbb E h_-}{2}.
\]

Expanding the covariance gives the exact identity \[
\operatorname{Cov}(g,h)
=\frac{\operatorname{Cov}(g_+,h_+)+\operatorname{Cov}(g_-,h_-)}2+d_gd_h.
\tag{2}
\]

Indeed, the joint expectation is the average of the two restricted joint
expectations, while the product of the two full means differs from the
average of the two restricted mean-products by \(d_gd_h\).

Since \(g_+,g_-\) take values in \(\{-1,1\}\), their half-difference is
in \(\{-1,0,1\}\). Therefore \[
|d_g|
\leq\mathbb E\left|\frac{g_+-g_-}{2}\right|
=\Pr[g_+\ne g_-]
=I_n(g).
\tag{3}
\]

Also \(|d_g|\leq1\). The same statements hold for \(h\). Hence \[
|d_gd_h|\leq\min\{I_n(g),I_n(h)\}.
\tag{4}
\]

Apply the triangle inequality and the inductive hypothesis to (2). For
every \(i<n\), \[
I_i(g)=\frac{I_i(g_+)+I_i(g_-)}2,
\qquad
I_i(h)=\frac{I_i(h_+)+I_i(h_-)}2.
\tag{5}
\]

For any nonnegative \(a,b,c,d\), \[
\frac{\min(a,b)+\min(c,d)}2
\leq\min\left(\frac{a+c}2,\frac{b+d}2\right),
\tag{6}
\]

because the left side is at most both terms inside the minimum. Apply
(6) to the two restricted influence pairs, use (5), and add (4). This
proves (1). \(\square\)

\section{The logarithmic split
inequality}\label{the-logarithmic-split-inequality}

\subsection*{Lemma 2}\label{lemma-2}
\addcontentsline{toc}{subsection}{Lemma 2}

Let \(a,b\) be nonnegative integers and \(A,B\geq0\) real numbers.
Suppose \(a=0\) implies \(A=0\), and \(b=0\) implies \(B=0\). Then \[
2\min(A,B)+A w(a)+B w(b)
\leq(A+B)w(a+b).
\tag{7}
\]

\textbf{Proof.} If \(a=0\), then \(A=0\), and both sides reduce to
\(B w(b)\). The case \(b=0\) is symmetric.

Suppose \(a,b>0\). Define \[
u=\log_2\frac{a+b}{a},\qquad
v=\log_2\frac{a+b}{b}.
\]

Both are nonnegative, and \[
u+v=\log_2\frac{(a+b)^2}{ab}\geq2,
\]

because \((a+b)^2-4ab=(a-b)^2\geq0\). Thus \[
Au+Bv\geq\min(A,B)(u+v)\geq2\min(A,B).
\]

Expanding \(u,v\) gives (7). \(\square\)

For a Lean implementation using natural logarithms, the nonzero-count
form is equivalently \[
2\ln(2)\min(A,B)+A\ln a+B\ln b
\leq(A+B)\ln(a+b).
\]

\section{Weighted tree theorem}\label{weighted-tree-theorem}

\subsection*{Theorem 3}\label{theorem-3}
\addcontentsline{toc}{subsection}{Theorem 3}

Every Boolean decision tree satisfying the no-repeat-on-a-path
convention obeys \[
\mathcal A(T)\leq P(T)=\sum_i I_i(f)w(k_i(T)).
\tag{8}
\]

Consequently, \[
\mathcal C(T)\leq\mathcal A(T)\leq\log_2(k)I(f)
\tag{9}
\]

for every read-\(k\) tree and \(k\geq1\).

\textbf{Proof.} Use structural induction on \(T\). For a leaf, every
influence and both covariance quantities are zero.

Suppose the root queries \(j\), and its subtrees \(T_+,T_-\) compute
\(g,h\). Since no path repeats \(j\), neither subtree queries \(j\).
Thus \(k_j(T)=1\) and \(I_j(g)=I_j(h)=0\). The root coordinate
contributes zero to all three potentials because \(w(1)=0\).

For \(i\ne j\), write \[
a_i=k_i(T_+),\quad b_i=k_i(T_-),\quad
A_i=I_i(g),\quad B_i=I_i(h).
\]

Then \[
k_i(T)=a_i+b_i,\qquad
I_i(f)=\frac{A_i+B_i}{2}.
\tag{10}
\]

If \(a_i=0\), then \(g\) does not depend on \(i\), so \(A_i=0\). The
analogous statement holds for \(b_i,B_i\). Lemma 2 therefore applies,
including every zero-count case. Summing it gives \[
\begin{aligned}
2P(T)-P(T_+)-P(T_-)
&=\sum_{i\ne j}\left[(A_i+B_i)w(a_i+b_i)-A_iw(a_i)-B_iw(b_i)\right]\\
&\geq2\sum_{i\ne j}\min(A_i,B_i)\\
&\geq2|\operatorname{Cov}(g,h)|,
\end{aligned}
\tag{11}
\]

where the last inequality is Lemma 1; coordinate \(j\) has zero
influence on both functions. By the induction hypothesis, \[
\begin{aligned}
\mathcal A(T)
&=|\operatorname{Cov}(g,h)|+\frac{\mathcal A(T_+)+\mathcal A(T_-)}2\\
&\leq|\operatorname{Cov}(g,h)|+\frac{P(T_+)+P(T_-)}2\\
&\leq P(T).
\end{aligned}
\]

This proves (8). If \(k_i(T)\leq k\), then \(w(k_i(T))\leq\log_2k\).
Multiply by the nonnegative influences and sum to get (9). \(\square\)

The signed inequality is all that is needed by the published Fourier
coding argument. The absolute version follows because the covariance
lemma already contains an absolute value and the potential proof has no
cancellations.

\section{Sharpness of the covariance
factor}\label{sharpness-of-the-covariance-factor}

The following family is the dummy-prefix example of Wan, Wright and Wu,
Section 2.4, specialized to a dictator. It is included to verify the
exact constant; the construction itself is prior work.

\subsection*{Proposition 4}\label{proposition-4}
\addcontentsline{toc}{subsection}{Proposition 4}

For every integer \(d\geq0\), there is a read-\(2^d\) decision tree
\(T_d\) computing a dictator and satisfying \[
I(f_d)=\operatorname{Var}(f_d)=1,\qquad
\mathcal C(T_d)=\mathcal A(T_d)=d,\qquad
P(T_d)=d.
\tag{12}
\]

\textbf{Proof.} Use coordinates \(z_1,\ldots,z_d,x\). First query
\(z_1,\ldots,z_d\) in a complete binary tree, independently of the
answers. At each of the resulting \(2^d\) nodes query \(x\) and output
its sign. The function is always \(f_d=x\). Thus \(I_x=1\), all
\(I_{z_r}=0\), and its variance is one.

The coordinate \(x\) occurs at \(2^d\) nodes. Coordinate \(z_r\) occurs
at \(2^{r-1}\) nodes, so the maximum count is \(2^d\). Every dummy-query
node has two identical child functions \(x\), hence covariance one. At
each depth \(r=0,\ldots,d-1\), there are \(2^r\) such nodes, each with
weight \(2^{-r}\), so that depth contributes one. The \(x\)-query nodes
have constant children and contribute zero. The total covariance is
therefore \(d\), and the weighted potential is \(I_x\log_2(2^d)=d\).
\(\square\)

This proves that the coefficient of \(I(f)\log_2k\) cannot be reduced
below one uniformly over all such trees. It does \textbf{not} prove that
logarithmic dependence is necessary in a Fourier entropy bound: a
dictator has Fourier entropy zero. Representation-dependent covariance
and function-dependent Fourier entropy must not be conflated.

\section{A four-variable counterexample to the proposed variance
bound}\label{a-four-variable-counterexample-to-the-proposed-variance-bound}

\subsection*{Proposition 5}\label{proposition-5}
\addcontentsline{toc}{subsection}{Proposition 5}

There is a read-three Boolean decision tree \(T\) on four coordinates,
with no repeated coordinate on a root-to-leaf path, computing the
two-variable AND function \(f\), such that \[
\mathcal C(T)=\mathcal A(T)=\frac54
>\frac34\log_2 3
=\operatorname{Var}(f)\log_2 3.
\tag{13}
\]

\textbf{Construction.} Use coordinates \(x,y,z,t\), each taking values
in \(\{-1,1\}\). Let \(D_x\) be the one-query dictator tree labelled
\(x\), with leaves \(+1\) and \(-1\) in its positive and negative
branches, respectively. Define \(D_y\) analogously. Write
\(N_i(S_-,S_+)\) for the tree with root query \(i\), negative subtree
\(S_-\) first, and positive subtree \(S_+\) second. With \(-\) denoting
the leaf labelled \(-1\), put \[
\begin{aligned}
L&=N_x\bigl(-,N_t(D_y,D_y)\bigr),\\
R&=N_y\bigl(-,N_t(D_x,D_x)\bigr),\\
T&=N_z(L,R).
\end{aligned}
\tag{14}
\]

Thus the negative branch of the root queries \(x\); after a positive
answer it asks the irrelevant coordinate \(t\), then asks \(y\),
independently of the answer to \(t\). The positive branch does the same
with \(x\) and \(y\) interchanged. A negative answer to the first
substantive query leads directly to the negative leaf. This branch
convention agrees with the Lean witness, whose first child is the
false/negative branch.

\textbf{Proof.} Each of \(L,R,T\) computes \[
f(x,y,z,t)=
\begin{cases}
+1,&x=y=+1,\\
-1,&\text{otherwise}.
\end{cases}
\tag{15}
\]

The two possible full query sequences are \((z,x,t,y)\) and
\((z,y,t,x)\), together with their shorter terminating prefixes.
Therefore no coordinate is repeated along a path. The counts are \[
k_x(T)=3,\qquad k_y(T)=3,\qquad
k_t(T)=2,\qquad k_z(T)=1.
\tag{16}
\]

The tree has nine internal nodes, and its read number is exactly three.
A uniform input gives \(f=+1\) with probability \(1/4\), so \[
\mathbb E f=-\frac12,\qquad
\operatorname{Var}(f)=1-\left(-\frac12\right)^2=\frac34.
\tag{17}
\]

Both root child functions equal \(f\), so their covariance is \(3/4\).
At the root of \(L\), one child computes the dictator \(y\) and the
other is constant, hence that node contributes zero covariance. The
dummy \(t\)-query has two dictator children, whose covariance is one.
Since this node is at depth one within \(L\), it follows that
\(\mathcal C(L)=1/2\). Its descendant dictator nodes have constant
children and hence zero covariance. The calculation for \(R\) is
identical. Consequently, \[
\mathcal C(T)=\frac34+\frac12\left(\frac12+\frac12\right)=\frac54.
\tag{18}
\]

All node covariances just computed are nonnegative, so
\(\mathcal A(T)=\mathcal C(T)\). Finally, \[
\log_2 3<\frac53
\quad\Longleftrightarrow\quad
3\log_2 3<5
\quad\Longleftrightarrow\quad
27<32,
\]

which is true. Multiplication by \(3/4\) proves (13). \(\square\)

\subsection{Scope of the refutation}\label{scope-of-the-refutation}

The proposal appears immediately after Theorem 2.8 in Section 2.4 of
Wan, Wright and Wu (2013/2014), and is explicitly called a conjecture in
Section 6.2 of Shalev (2018/2019). It concerns the covariance of the
given read-\(k\) representation. The original paper's decision-tree
convention in Section 2.1 rules out repeated queries along a path; it
does not require the tree to have minimum size, to have minimum expected
depth, or to omit irrelevant queries. Its Figure 2 and accompanying
discussion explicitly use dummy-query layers with identical child
functions. The construction (14) respects precisely that scope.

This counterexample does not address an existential statement asserting
that every function has \emph{some} suitably efficient representation
obeying a logarithmic variance estimate. The AND function itself has a
read-once tree with zero tree covariance. The refuted statement
quantifies over all given read-\(k\) trees, as in the proposed
improvement to the published covariance lemma.

\section{Counterexample with every query
relevant}\label{counterexample-with-every-query-relevant}

\subsection*{Proposition 6}\label{proposition-6}
\addcontentsline{toc}{subsection}{Proposition 6}

There is a nine-variable read-three tree, with no coordinate repeated on
a path and with different child functions at every internal node, for
which \[
\mathcal C(T)=\frac{4863}{4096},\qquad
\operatorname{Var}(f)=\frac{3007}{4096},\qquad
\mathcal C(T)>\log_2(3)\operatorname{Var}(f).
\tag{19}
\]

\textbf{Construction.} Use Boolean coordinates
\(z,x,y,u,v,r_1,r_2,r_3,r_4\), with false and true encoded as -1 and +1
for computing moments. Let \(R=r_1\lor r_2\lor r_3\lor r_4\),
implemented by its short-circuit OR tree. Use \(Q_i\) for a single query
of coordinate \(i\), and write \(N_i(A,B)\) with the false branch first
and the true branch second, as above. Put \[
\begin{aligned}
L&=N_x\bigl(-,N_u(Q_y,N_y(-,R))\bigr),\\
M&=N_y\bigl(-,N_v(Q_x,N_x(-,R))\bigr),\\
T&=N_z(L,M).
\end{aligned}
\tag{20}
\]

Thus the left child function is the signed encoding of
\(x\land y\land(\neg u\lor R)\), while the right child function is the
signed encoding of \(x\land y\land(\neg v\lor R)\). The same four OR
coordinates are used in both branches. The exact query-count vector, in
the displayed coordinate order, is \[
(1,3,3,1,1,2,2,2,2).
\tag{21}
\]

No path repeats a coordinate. Each OR query has distinct child
functions. The \(x\) and \(y\) queries gate a nonconstant positive
branch; the \(u\) and \(v\) queries distinguish a dictator from its
conjunction with \(R\). At the root, taking \(x=y=1\), all \(r_i=0\),
and \(u\ne v\) distinguishes the two child functions. Hence every
internal query affects its subtree's output. The Lean theorem checks
inequality of the actual child functions, not merely syntactic
inequality of the subtrees.

\textbf{Exact moment calculation.} Put \(\epsilon=\Pr[R=0]=1/16\). Both
child functions \(g,h\) have positive-output probability \[
p=\frac14\left(1-\frac\epsilon2\right)=\frac{31}{128},
\qquad \mathbb E g=\mathbb E h=2p-1=-\frac{33}{64}.
\]

Since the root bit is independent, \(\mathbb E f=-33/64\), and therefore
\(\operatorname{Var}(f)=1-(33/64)^2=3007/4096\). The probability that
both child outputs are positive is \[
p_{++}=\frac14\left((1-\epsilon)+\frac\epsilon4\right)=\frac{61}{256}.
\]

Consequently \[
\mathbb E[gh]=1-4p+4p_{++}=\frac{63}{64},\qquad
\operatorname{Cov}(g,h)=\frac{2943}{4096}.
\]

For a dictator \(Y\) and the signed encoding of \(Y\land R\), the
covariance is \(1-\epsilon=15/16\): the dictator has mean zero, and
their product is negative exactly when \(Y=1\) and \(R=0\). The outer
query of each child tree has a constant branch and hence covariance
zero. Its selector query contributes \(15/16\) at depth one; all
remaining nodes also have a constant branch and contribute zero. Thus \[
\mathcal C(L)=\mathcal C(M)=\frac{15}{32},\qquad
\mathcal C(T)=\frac{2943}{4096}+\frac{15}{32}=\frac{4863}{4096}.
\]

Finally, \(3^5=243<256=2^8\) gives \(\log_2 3<8/5\), and \[
8\cdot3007=24056<24315=5\cdot4863.
\]

Multiplying these strict inequalities by the appropriate positive
denominators proves (19). This counterexample satisfies the additional
local relevance condition as well as the original tree convention. It is
not asserted to be a tree of minimum size or minimum expected depth.
\(\square\)

\section{Relation to previous work}\label{relation-to-previous-work}

Wan, Wright and Wu introduced tree covariance to bound the expected
length of a protocol for encoding a function's Fourier spectrum. Their
Lemmas 1.5 and 1.7 connect this protocol to Fourier entropy. Their
covariance estimate is \(\mathcal C(T)\leq(k-1)\operatorname{Var}(f)\).
Section 2.4 of their paper proposes the stronger bound
\(\mathcal C(T)\leq\log_2(k)\operatorname{Var}(f)\) and gives the
dummy-prefix equality construction used in Proposition 4. Propositions 5
and 6 refute that stronger bound for a given tree.

Shalev's Section 6.2, Theorem 31, proves
\(\mathcal C(T)\leq2\sqrt{k}I(f)\). The same section proposes the
Fourier weight \(\ell_T(S)=\sum_{i\in S}\log_2 k_i(T)\). After
interchanging finite sums, this gives the potential in Theorem 3, since
\(I_i(f)=\sum_{S\ni i}\widehat f(S)^2\). The logarithmic potential is
therefore due to Shalev. Inequality (1) proves its covariance bound
through an induction on the tree. The theorem also bounds the sum of
absolute node covariances.

\subsection{Distinct read parameters and selector
hypotheses}\label{distinct-read-parameters-and-selector-hypotheses}

Arunachalam et al.~prove a logarithmic read-number bound for Fourier
min-entropy of disjunctive normal form (DNF) formulas. A DNF is a
disjunction of conjunctions of literals. Its read number counts terms
that contain a variable. Shmalo's Theorem 5.1 gives Fourier-entropy
bounds for unions and intersections of structured gates. Its read
parameter counts constituent gates that depend on a coordinate. These
parameters differ from the query-node count in a decision tree.

In a tree, \(k_i\) counts query nodes. Converting each accepting
root-to-leaf path into a DNF term can repeat a root variable in
arbitrarily many terms, even when the tree queries it only once. This
conversion does not preserve the read number \(k\).

Shmalo's Proposition 8.6 gives general identities for selectors with
shared data. Those identities do not require balanced data functions or
disjoint Fourier supports. Corollaries 8.7 and 8.8 impose further
hypotheses, including orthogonality conditions between neighboring
branches. Arbitrary tree children need not satisfy these conditions. For
example, both children can be the same dictator and thus have the same
nonempty Fourier support. The general identities alone do not give the
covariance estimate (8).

González, MacManus and Pereyra obtain Fourier entropy-influence bounds
for further function classes through a restriction inequality. Han
bounds Fourier entropy in terms of total influence and coordinate
influences with reciprocal-logarithm weights. These results concern
Fourier entropy; neither is a statement of the representation-dependent
covariance bound proved here.

\subsection{Consequences and limits}\label{consequences-and-limits}

Theorem 3 replaces the square-root dependence in Shalev's covariance
estimate by a logarithm. Proposition 4 shows that its coefficient one is
optimal for infinitely many read numbers. The example computes a
dictator, whose Fourier entropy is zero. It therefore establishes
sharpness of the covariance estimate, not sharpness of a Fourier-entropy
bound.

The logarithmic influence estimate and the earlier variance estimate
give different information. Influence can be much larger than variance,
so Theorem 3 need not improve the earlier variance estimate for every
function. Propositions 5 and 6 refute the coefficient-one logarithmic
variance proposal. They do not determine the optimal asymptotic bound in
terms of variance and read number alone.

The formal development contains no Fourier transform, Shannon entropy,
or coding theorem. Its results are the covariance inequalities, their
equality family, and the finite counterexamples. The cited encoding
protocol gives the motivation. The unrestricted Fourier
entropy-influence conjecture is not resolved here.

\section{Lean verification and
reproducibility}\label{lean-verification-and-reproducibility}

The accompanying project formalizes the results proved in this note in
Lean 4.19.0, using mathlib revision
\texttt{c44e0c8ee63ca166450922a373c7409c5d26b00b}. The project defines
its own finite Boolean cube, uniform expectation, bit-flip influence,
tree syntax, evaluation, global label counts, and signed and absolute
tree covariance. The regularity predicate rules out repeated coordinates
on a root-to-leaf path. No theorem is replaced by a hypothesis asserting
the desired conclusion.

The main correspondence is:

{\def\LTcaptype{none} % do not increment counter
\begin{longtable}[]{@{}
  >{\raggedright\arraybackslash}p{(\linewidth - 2\tabcolsep) * \real{0.1978}}
  >{\raggedright\arraybackslash}p{(\linewidth - 2\tabcolsep) * \real{0.8022}}@{}}
\toprule\noalign{}
\begin{minipage}[b]{\linewidth}\raggedright
Result
\end{minipage} & \begin{minipage}[b]{\linewidth}\raggedright
Formal declaration, within the \texttt{\seqsplit{ReadK}} namespace
\end{minipage} \\
\midrule\noalign{}
\endhead
\bottomrule\noalign{}
\endlastfoot
Lemma 1 & \texttt{\seqsplit{abs\_cov\_le\_sum\_min}} \\
Lemma 2 & \texttt{\seqsplit{logCount\_potential}} \\
Theorem 3, weighted &
\texttt{\seqsplit{Tree.absoluteCovariance\_le\_potential}} \\
Theorem 3, read bound &
\texttt{\seqsplit{Tree.absoluteCovariance\_le\_log\_read\_mul\_influence}} \\
Proposition 4 & \texttt{\seqsplit{Sharpness.full\_equality\_family}} \\
Proposition 5 &
\texttt{\seqsplit{Counterexample.variance\_refinement\_false}} and the
exact witness statistics \\
Proposition 6 &
\texttt{\seqsplit{Nonredundant.violates\_log\_variance}}, with its
regularity, relevance, count, and moment theorems \\
\end{longtable}
}

The development uses no proof holes or new axiom declarations. The
transitive axiom audit is included in \texttt{verification/axioms.txt};
its permitted foundational axioms are propositional extensionality,
classical choice, and quotient soundness. Finite decision proofs and
algebraic tactics produce terms checked by Lean's kernel. No native
evaluation oracle is used.

\newpage

After installing Lean's standard toolchain manager, reproduce the
project from its root directory with:

\begin{Shaded}
\begin{Highlighting}[]
\ExtensionTok{lake}\NormalTok{ exe cache get}
\FunctionTok{bash}\NormalTok{ verify.sh}
\ExtensionTok{python3}\NormalTok{ scripts/test\_verifiers.py}
\end{Highlighting}
\end{Shaded}

The toolchain file and Lake manifest pin the compiler and dependencies.
Keep the supplied manifest when reproducing these results. The optional
cache step downloads official compiled mathlib dependencies; the
verification script rebuilds this project's proof sources, prints the
theorem types and their axiom dependencies, scans the sources, and runs
an independent exact-rational check of both finite counterexamples. The
Python computation is a supplementary check and does not replace any
Lean proof. Omitting the cache step causes dependencies to be built from
source.

The verifier tests check that invalid or incomplete evidence produces a
nonzero exit status. The package includes a summary of the cited prior
work in \texttt{docs/related-work.md}. Formal verification establishes
the statements in the specified finite model. It does not establish
literature priority.

\section*{Authorship and assistance}\label{authorship-and-assistance}
\addcontentsline{toc}{section}{Authorship and assistance}

Dev D. Goyal directed the research and is responsible for publication.
ChatGPT (OpenAI) contributed mathematical arguments and Lean
formalizations. Codex assisted with verification and editing.

\section*{References}\label{references}
\addcontentsline{toc}{section}{References}

\begin{enumerate}
\def\labelenumi{\arabic{enumi}.}
\tightlist
\item
  Andrew Wan, John Wright, and Chenggang Wu. \emph{Decision Trees,
  Protocols, and the Fourier Entropy-Influence Conjecture}. ITCS 2014,
  pp.~67--80. DOI: \url{https://doi.org/10.1145/2554797.2554806}.
  Preprint: \url{https://arxiv.org/abs/1312.3003}.
\item
  Guy Shalev. \emph{On the Fourier Entropy Influence Conjecture for
  Extremal Classes}. arXiv:1806.03646, v2, 24 January 2019.
  \url{https://arxiv.org/abs/1806.03646v2}.
\item
  Srinivasan Arunachalam, Sourav Chakraborty, Michal Koucký, Nitin
  Saurabh, and Ronald de Wolf. \emph{Improved Bounds on Fourier Entropy
  and Min-entropy}. ACM Transactions on Computation Theory 13(4),
  Article 22, 2021. DOI: \url{https://doi.org/10.1145/3470860}.
  Preprint: \url{https://arxiv.org/abs/1809.09819}.
\item
  María José González, Paul MacManus, and María Cristina Pereyra.
  \emph{Further evidence towards the Fourier Entropy-Influence
  conjecture}. arXiv:2606.00246v2, 9 June 2026.
  \url{https://arxiv.org/abs/2606.00246v2}.
\item
  Yitzchak Shmalo. \emph{Fourier entropy of threshold functions and
  their compositions}. arXiv:2610.16069v1, 12 October 2026.
  \url{https://arxiv.org/abs/2610.16069v1}.
\item
  Xiao Han. \emph{A New Bound for the Fourier-Entropy-Influence
  Conjecture}. Combinatorica 45, Article 4, 2025.
  \url{https://doi.org/10.1007/s00493-024-00133-z}.
\end{enumerate}

\end{document}
